\chapter{Discussion}
\label{ch:discussion}

\section{Interpretation of Main Results}
\label{sec:interpretation}

The cohort wage gap estimates reveal meaningful individual-level returns to ProUni that are masked when pooling across age groups. The headline finding --- a positive and significant ATT on the young--old wage gap in low-dose microregions --- indicates that ProUni raised the relative earnings of workers in the age range plausibly exposed to the program. A back-of-envelope calculation using the Wald/LATE estimate implies a return of approximately [FILL]\% per ProUni-induced degree, which [FILL: is consistent with / exceeds / falls below] the Mincerian estimates of 15--20\% returns to higher education in Brazil documented by \citeA{BarbosaFilhoPessoa2008}. The old-cohort placebo confirms that this effect is age-specific and cannot be attributed to general local prosperity associated with ProUni rollout.

\section{Scarcity vs.\ Saturation}
\label{sec:saturation}

The finding that low-dose microregions exhibit larger wage returns than high-dose microregions is consistent with the diminishing returns hypothesis articulated by \citeA{Duflo2001} in the context of Indonesia's school construction program. Where ProUni scholarships are scarce, each additional degree holder captures a large college wage premium. Where scholarships are abundant, the supply of college-educated workers increases faster than demand, compressing the premium through general equilibrium effects. This saturation pattern has important implications for the optimal scale of higher-education subsidies: ProUni's returns are highest precisely where expansion is most constrained, suggesting that further geographic targeting could improve cost-effectiveness.

\section{Heterogeneity: Gender and Race}
\label{sec:hetero_discussion}

The gender heterogeneity results [FILL: describe pattern]. This connects to ProUni's affirmative action mandate, which reserves scholarships for students from public secondary schools and low-income families, with racial and disability quotas. The fact that [FILL: female / non-white] workers show larger treatment effects suggests that the program successfully reaches populations with higher marginal returns to education, potentially narrowing longstanding labor market inequalities. This finding complements the glass ceiling literature in Brazil, which documents persistent gender and racial wage gaps even conditional on education \shortciteA{[FILL: relevant citation]}.

\section{Geographic Concentration}
\label{sec:geography}

The geographic heatmap (Figure~\ref{fig:map}) reveals stark regional patterns in ProUni intensity. High-dose microregions are concentrated in the Southeast (S\~{a}o Paulo, Minas Gerais, Rio de Janeiro), where private universities are most abundant. Low-dose microregions predominate in the North and Northeast, where the private higher-education sector is thinner. This geographic distribution implies that ProUni's wage effects are strongest in precisely the regions that have historically received less investment in human capital, suggesting a potential role for the program in reducing regional inequality --- albeit with diminishing returns at higher doses.

\section{Policy Implications}
\label{sec:policy}

Three policy implications emerge from the analysis. First, ProUni works: it generates measurable wage returns for its beneficiaries, validating the program's core design of leveraging private-sector capacity for public human capital investment. Second, returns diminish at scale, implying that marginal expansions in already-saturated markets (e.g., S\~{a}o Paulo) may yield lower returns than redirecting resources to underserved regions. Third, the complementarity between ProUni and other policies --- FIES (student loans), Lei de Cotas (affirmative action quotas at public universities), and state-level programs --- deserves further investigation, as their combined effects likely differ from the sum of their individual impacts.

\section{Limitations}
\label{sec:limitations}

Several limitations warrant discussion. First, the analysis uses formal-sector data from RAIS, which excludes the informal sector (approximately 40\% of Brazilian employment). If ProUni graduates disproportionately enter formal employment, the wage estimates may overstate the population-level effect. Second, the ecological design cannot track individual scholarship recipients to their labor market outcomes; the cohort decomposition mitigates but does not eliminate this limitation. Third, the treatment variable (cumulative scholarships per 1,000 youth) measures program intensity rather than actual degree completion, introducing measurement error that likely attenuates the estimates. Fourth, the simultaneous expansion of FIES during the analysis period represents a potential confound, though the Duflo-style age decomposition partially addresses this by exploiting age-specific rather than region-specific variation. Fifth, the 22--28 age window for the exposed cohort is an approximation; some workers in this range may not be ProUni beneficiaries, and some older workers may have received scholarships through non-standard pathways.
